\begin{textsample}{NEG Dim 1 – 2015\_c4 – Score -4.00 – 2015\_c4/enade\_2015\_tec\_gest\_qual\_d7c5047481.txt}  \label{ex:f1_neg_106}
CAPÍTULO 4 PERCEPÇÃO DA PROVA As análises feitas neste capítulo tratam das percepções dos concluintes da Área de Tecnologia em Gestão da Qualidade sobre a prova aplicada no Enade/2015.
Essas percepções foram mensuradas por meio de nove questões que avaliaram desde o grau de dificuldade da prova até o tempo gasto para concluí -la.
As percepções sobre a prova foram relacionadas com o desempenho dos estudantes e com a Grande Região de funcionamento do curso.
O questionário de percepção da prova encontra -se ao final do Anexo VII, que traz a reprodução da prova.
O desempenho dos estudantes foi classificado em quatro quartos.
Para tanto, esse desempenho foi ordenado de forma ascendente.
O percentil 25, P25, também conhecido como primeiro quartil, é a nota de desempenho que deixa um quarto ( 25 \% ) dos valores observados abaixo e três quartos acima.
A Figura 1 apresenta uma ilustração deste conceito.
O quarto inferior de desempenho é composto pelas notas abaixo do primeiro quartil.
Já o percentil 75, P75, também conhecido como terceiro quartil, é o valor para o qual há três quartos ( 75 \% ) dos dados abaixo e um quarto acima dele.
O quarto superior de desempenho é composto pelas notas iguais ou acima do terceiro quartil.
O percentil 50, P50, também conhecido como mediana, é o valor que divide as notas em dois conjuntos de igual tamanho.
O segundo quarto inclui valores entre o primeiro quartil ( P25 ) e a mediana.
O terceiro quarto contém os valores entre a mediana ( P50 ) e o terceiro quartil ( P75 ).
Vale ressaltar que percentis, quartis e medianas são pontos que não obrigatoriamente pertencem ao conjunto original de dados, ao passo que os quartos são subconjuntos dos dados originais.
A seguir, serão apresentados gráficos com resultados selecionados, relativos às nove questões avaliadas por grupos de estudantes.
Os gráficos apresentam nas barras o percentual de alunos que assinalaram uma das opções ou a soma das porcentagens daqueles que assinalaram duas ( ou três ) delas.
Por exemplo, para as questões 1 e 2, os gráficos apresentam a porcentagem total de participantes que assinalaram as opções ( D ) Difícil e ( E ) Muito difícil.
Em \textbf{cada} barra foram assinalados também os extremos do intervalo de confiança de 95 \% como linhas verticais unidas por uma linha horizontal na forma da letra H maiúscula.
As Tabelas no Anexo II apresentam os valores absolutos e a distribuição percentual¹⁷ das alternativas válidas das nove questões, segundo o mesmo recorte de desempenho dos alunos e Grande Região de funcionamento do curso. 4.1 GRAU DE DIFICULDADE DA PROVA 4.1.1 Componente de Formação Geral Ao avaliarem “ Qual o grau de dificuldade desta prova na parte de Formação Geral? ” ( Questão 1 ), 36,0 \% do grupo de inscritos e presentes optaram pelas alternativas Difícil ou Muito difícil.
Entretanto, para mais da metade dos estudantes ( 57,9 \% ), o Componente de Formação Geral da prova foi considerado com grau de dificuldade Médio ( Gráfico 4.1, Gráfico 4.2 e, no Anexo II, a Tabela II.1 ).
O percentual de estudantes que consideraram a prova como Difícil ou Muito difícil foi maior na região Sul, onde a proporção foi de 38,5 \%, enquanto a de menor incidência foi a Nordeste, com 21,1 \%.
No Gráfico 4.1, é possível observar diferenças estatisticamente significativas entre a região Nordeste e Sudeste ( 37,7 \% ) e Sul.
Nas Grandes Regiões, a proporção de presentes à prova que consideraram o Componente de Formação Geral como sendo de grau de dificuldade Médio \textbf{esteve} entre 54,8 \% na região Sul e 71,8 \% na região Nordeste. ¹⁷ Cumpre \textbf{lembrar} uma das convenções para tabelas numéricas ( pág. iii ) sobre a possibilidade de a soma das partes não resultar em 100 \% por questões de arredondamento.
O percentual de alunos que consideraram a prova Difícil ou Muito difícil foi 40,2 \% no primeiro quarto e 34,9 \% no quarto, grupo de melhor desempenho na prova.
Nos quartos de desempenho intermediários, a proporção de alunos que consideraram a prova Difícil ou Muito difícil foi de 35,8 \% no segundo quarto e 33,0 \% no terceiro quarto.
Não se observam diferenças estatisticamente significativas entre os quartos.
Para todos os quartos de desempenho, a alternativa modal para esta pergunta foi Médio, com 54,9 \% e 58,1 \% dos respondentes nos quartos extremos, primeiro e quarto, respectivamente, sendo a maior proporção 60,2 no terceiro quarto. 4.1.2 Componente de Conhecimento Específico Ao responderem à Questão 2 – “ Qual o grau de dificuldade desta prova na parte de Componente Específico? ” – 40,5 \% do grupo de estudantes classificaram -na como Difícil ou Muito difícil.
Além disso, o Componente de Conhecimento Específico da prova foi considerado com grau de dificuldade Médio por 54,5 \% dos alunos ( Gráfico 4.3, Gráfico 4.4, e, no Anexo II, a Tabela II.2 ).
A análise das respostas dos estudantes quanto ao grau de dificuldade do Componente de Conhecimento Específico da prova, agregado por Grande Região, mostra que a diferença entre a maior e as duas menores proporções de alunos que a avaliaram como Difícil ou Muito difícil é estatisticamente significativa : a maior na região Sul ( 44,3 \% ) e as menores proporções nas regiões Nordeste ( 16,9 \% ) e Norte ( 30,4 \% ).
Também são observadas diferenças estatisticamente significativas entre a região Sudeste ( 43,3 \% ) e as regiões Norte e Nordeste.
O percentual de alunos que classificou o grau de dificuldade como Médio, no Componente de Conhecimento Específico, variou de 51,0 \% a 74,6 \%, para as regiões Sul e Nordeste, respectivamente.
Considerando -se a avaliação da dificuldade das questões do Componente de Conhecimento Específico da prova, de acordo com o desempenho dos estudantes, a proporção dos que classificaram a parte específica como Difícil ou Muito difícil variou de 37,1 \% ( terceiro quarto ) a 43,8 \% ( segundo quarto ).
Nos quartos extremos dos que optaram por essa categoria, os percentuais foram 40,8 \% ( quarto inferior ) e 40,5 \% ( quarto superior ).
Não se observam diferenças estatisticamente significativas entre os resultados dos quatro quartos de desempenho.
A alternativa modal para a Questão 2 foi o grau Médio, variando de 52,2 \% ( segundo quarto ) a 58,2 ( terceiro quarto ).
Nos quartos extremos os percentuais foram 54,8 \% ( no inferior ) e 52,9 \% ( no superior ) 4.2 EXTENSÃO DA PROVA EM RELAÇÃO AO TEMPO TOTAL Indagados quanto à extensão da prova, em relação ao tempo total oferecido para a sua resolução ( Questão 3 ), os estudantes apontaram, com maior incidência, a alternativa que considerava a extensão adequada, para todas as agregações consideradas ( Gráfico 4.5, Gráfico 4.6, e, no Anexo II, a Tabela II.3 ).
O percentual de alunos que responderam ser a extensão da prova adequada foi de 60,2 \%.
Já 33,5 \% dos inscritos presentes consideraram que a prova foi longa ou muito longa e menos do que 7 \% a avaliaram como curta ou muito curta.
Entre as Grandes Regiões a proporção daqueles que avaliaram a prova como longa ou muito longa em relação ao tempo total destinado à sua resolução variou de 27,3 \% na região Norte até 36,5 \% na região Sul, não havendo diferenças estatisticamente significativas entre as regiões.
A alternativa modal para essa questão foi extensão adequada, com proporções variando de 54,9 \% ( Nordeste ) a 61,3 \% ( Sudeste ).
Considerando -se o desempenho dos alunos, nota -se ainda que 55,6 \% consideraram a extensão da prova adequada no quarto de desempenho inferior e 61,1 \% no de melhor desempenho ( quarto superior ).
Nos quartos intermediários, esta proporção foi 60,1 \% no segundo quarto e 63,9 \% no terceiro.
No Gráfico 4.6, pode -se constatar que a proporção de estudantes que consideraram a prova longa ou muito longa em função dos quartos de desempenho variou de 37,4 \% no primeiro quarto a 30,3 \% no terceiro.
No quarto superior a proporção foi 33,9 \%.
Não se observam diferenças estatisticamente significativas entre as proporções de escolha entre os quartos de desempenho. 4.3 COMPREENSÃO DOS ENUNCIADOS DAS QUESTÕES 4.3.1 Componente de Formação Geral Com relação aos enunciados das questões do Componente de Formação Geral ( Questão 4 ), as opiniões foram positivas, já que 73,2 \% dos alunos avaliados consideraram todos ou a maioria dos enunciados das questões “ claros e objetivos ” ( Gráfico 4.7, Gráfico 4.8, e, no Anexo II, a Tabela II.4 ).
Na análise regional, a percentagem de estudantes que avaliaram que todos ou a maioria dos enunciados das questões do Componente de Formação Geral \textbf{estavam} “ claros e objetivos ” variou de 69,4 \% na região Nordeste a 74,6 \% na região Norte, não havendo diferenças estatisticamente significativas entre as regiões.
A análise das percepções dos estudantes sobre a clareza e objetividade dos enunciados permite afirmar que todos, ou a maioria dos enunciados de questões relativas ao Componente de Formação Geral, foram considerados “ claros e objetivos ” para a maior parte dos respondentes ( maior ou igual a 69,4 \% em todas as regiões e maior ou igual a 65,5 \% para todos os quartos de desempenho ).
Segundo o desempenho, observa -se que a proporção dos que emitiram esta opinião cresce conforme o desempenho aumenta, com diferenças estatisticamente significativas do primeiro para o terceiro e o último quartos.
No quarto superior, a clareza e objetividade de todos ou da maioria dos enunciados das questões foi percebida por 77,4 \% dos alunos sem que seja estatisticamente significativa a diferença deste para o quarto imediatamente inferior, onde a proporção foi de 76,0 \%. 4.3.2 Componente de Conhecimento Específico Com relação aos enunciados das questões do Componente de Conhecimento Específico da prova, para 71,8 \% dos estudantes avaliados da Área de Tecnologia em Gestão da Qualidade, a clareza e a objetividade ( Questão 5 ) \textbf{estavam} presentes em todas ou na maioria das questões ( Gráfico 4.9, Gráfico 4.10, e no Anexo II, a Tabela II.5 ).
A maioria dos estudantes de todas as Grandes Regiões brasileiras considerou “ claros e objetivos ” todos ou a maioria dos enunciados das questões do Componente de Conhecimento Específico da prova, percentual sempre maior do que 70 \%.
Não há diferenças estatisticamente significativas entre as regiões.
A proporção de estudantes que consideraram os enunciados das questões “ claros e objetivos ” apresenta uma tendência crescente em relação ao aumento de desempenho : mais elevada no quarto superior ( 78,0 \% ) se comparada ao quarto inferior de desempenho ( 63,8 \% ).
As diferenças do quarto inferior para os dois últimos quartos são estatisticamente significativas. 4.4 SUFICIÊNCIA DAS INFORMAÇÕES/INSTRUÇÕES FORNECIDAS Ao avaliarem as informações/instruções fornecidas para a resolução das questões ( Questão 6 ), 79,0 \% dos respondentes da Área de Tecnologia em Gestão da Qualidade de todo o Brasil afirmaram que estas eram até excessivas ou “ suficientes ” em todas ou na maioria das questões ( Gráfico 4.11, Gráfico 4.12, e, no Anexo II, a Tabela II.6 ).
Quanto à distribuição de respondentes pelas Grandes Regiões observa -se que a proporção de estudantes que consideraram as informações/instruções fornecidas até excessivas ou “ suficientes ” em todas ou na maioria das questões foi sempre igual ou superior a 76,7 \%, chegando a 82,6 \% na região Sul.
O menor percentual foi o da região Norte ( 76,7 \% ).
Não há diferenças estatisticamente significativas entre as regiões.
Levando -se em conta o desempenho dos participantes, o percentual de participantes que avaliaram “ as informações/instruções ” como até excessivas ou “ suficientes ” em todas ou na maioria das questões foi mais elevado no quarto superior ( 85,8 \% ), percentual superior à média nacional ( 79,0 \% ).
Já no quarto inferior, a suficiência das informações/instruções declarada como até excessiva, em todas ou na maioria das questões foi percebida por 66,4 \% dos respondentes.
Estas proporções são crescentes com os quartos de desempenho.
Notam -se diferenças estatisticamente significativas entre as opiniões de estudantes do quarto inferior de desempenho em relação aos demais, assim como do segundo quarto em relação ao quarto de melhor desempenho, como mostra o Gráfico 4.12. 4.5 DIFICULDADE ENCONTRADA AO RESPONDER À PROVA \textbf{Perguntados} sobre as dificuldades com as quais se depararam ao responder à prova ( Questão 7 ), 11,9 \% dos estudantes apontaram o Desconhecimento do conteúdo.
Para 58,8 \%, a Forma diferente de abordagem do conteúdo foi indicada como dificuldade.
Já a Falta de motivação para fazer a prova foi a dificuldade apontada por 7,2 \% dos respondentes.
Considerando -se todo o Brasil, 18,1 \% dos respondentes afirmaram que não tiveram qualquer tipo de dificuldade para responder à prova ( Tabela II.7 no Anexo II ).
Os Gráficos 4.13 e 4.14 apresentam os percentuais de estudantes que apontaram o Desconhecimento do conteúdo como dificuldade percebida ao responder à prova.
Na análise por Grandes Regiões, o percentual de inscritos e presentes que apontaram o Desconhecimento do conteúdo como dificuldade ao responder à prova não superou 14 \%.
Os percentuais variaram de 9,0 \% na região Nordeste a 14,0 \% na Norte.
Não há diferenças estatisticamente significativas entre as regiões.
A Forma diferente de abordagem do conteúdo foi a escolha modal dos estudantes, com percentuais que variaram de 44,8 \% ( região Nordeste ) a 60,2 \% ( Sudeste ).
O percentual de alunos que citaram a Falta de motivação como dificuldade variou de 6,1 \% ( região Sudeste ) a 11,9 \% ( região Nordeste ).
Os que declararam não ter qualquer tipo de dificuldade para responder à prova variaram de 16,6 \% na região Sul a 25,4 \% na Nordeste.
Com relação aos quartos de desempenho, o Desconhecimento do conteúdo foi a opção escolhida por 12,3 \% dos estudantes do quarto superior e 12,2 \% do quarto inferior.
Não existem diferenças estatisticamente significativas entre os quartos de desempenho ( Gráfico 4.14 ).
A alternativa modal para os alunos, quando agregados pelos quartos de desempenho, foi que a dificuldade encontrada foi causada pela Forma diferente de abordagem do conteúdo : 61,9 \% no quarto inferior e 58,0 \% do quarto superior assim o responderam.
A proporção de alunos que citaram a Falta de motivação como dificuldade variou de 5,3 \% no quarto superior a 10,2 \% no inferior.
Os que declararam não ter qualquer tipo de dificuldade para responder à prova variaram de 13,2 \% no quarto inferior a 20,9 \% no superior. 4.6 CONTEÚDOS DAS QUESTÕES OBJETIVAS DA PROVA Ao analisarem os conteúdos das questões objetivas da prova ( Questão 8 ), um percentual muito pequeno dos estudantes avaliados, apenas 4,1 \%, afirmou que não estudou ainda a maioria desses conteúdos ( Gráficos 4.15, Gráfico 4.16, e a Tabelas II.8 no Anexo II ).
A maioria ( 66,3 \% ) afirmou ter estudado e aprendido muitos ou todos os conteúdos.
Na análise por Grande Região, a proporção de respondentes que escolheram a opção não estudou ainda a maioria desses conteúdos foi pequena.
Observa -se que nas regiões Norte ( 5,2 \% ) e Sudeste ( 4,6 \% ), apesar de pequenas, as proporções foram maiores do que a média nacional ( 4,1 \% ).
Não se observam diferenças estatisticamente significativas entre as regiões.
Em todas as regiões, a maioria dos presentes afirmou ter estudado e aprendido muitos ou todos os conteúdos, com proporções variando entre 60,8 \% na região Norte e 73,9 \% na Nordeste.
Considerando -se separadamente as opiniões de estudantes dos quatro quartos de desempenho, observa -se que, no quarto inferior, 7,2 \% ofereceram como resposta que não estudou ainda a maioria desses conteúdos, sendo 2,0 \% os do quarto superior com a mesma resposta.
A diferença entre os alunos que optaram por este motivo de dificuldade é estatisticamente significativa entre o primeiro e o terceiro e o último quartos de desempenho.
Tendo em conta o quarto superior, 79,8 \% dos alunos afirmaram ter estudado e aprendido muitos ou todos os conteúdos.
No outro extremo, no primeiro quarto, 52,3 \% optaram pelas mesmas categorias. 4.7 TEMPO GASTO PARA CONCLUIR A PROVA Ao responderem sobre o tempo de conclusão da prova ( Questão 9 ), mais de três quartos dos estudantes ( 81,0 \% ) afirmou ter gasto Entre duas e quatro horas ( Gráfico 4.17, Gráfico 4.18 e, no Anexo II, a Tabela II.9 ).
Considerando -se as cinco Grandes Regiões brasileiras, o percentual dos que utilizaram Entre duas e quatro horas para finalizar a prova variou de 76,9 \% na região Norte a 82,0 \% na região Sudeste, como mostra o Gráfico 4.17.
Não existem diferenças estatisticamente significativas entre as regiões.
Uma vez sendo analisadas as alternativas escolhidas pelos estudantes que se situam nos diferentes quartos de desempenho, observa -se uma tendência crescente : uma maior proporção de participantes no quarto superior declarou ter gasto Entre duas e quatro horas “ para concluir a prova ” quando comparadas com os dos quartos inferiores.
A diferença entre o primeiro ( 76,7 \% ) e o quarto de maior desempenho ( 84,7 \% ) é estatisticamente significativa.

% matched lemmas: cada, estar, lembrar, perguntar
\end{textsample}
